
\section{Discussion}
\label{sec:discussion}
% Our work on \tool introduces a significant paradigm shift, moving from standalone automated solvers to a collaborative, AI-guided reasoning process. The success of the \tool framework demonstrates that augmenting symbolic engines with learned, high-level policies is a highly promising direction. In this section, we reflect on the broader implications of this work, analyze its current limitations, and outline promising directions for creating even more powerful and intelligent reasoning partners.

\subsection{Limitations}
% While \tool demonstrates notable solving time reductions on structured constraints, its robustness varies. The RL agent’s effectiveness depends on the representativeness of the training set; unexpected constraint types may lead to poor variable selection. Furthermore, the LLM’s output quality can vary, and it may sometimes propose incorrect values. We mitigate this by having the solver verify each LLM-suggested assignment, but those extra solver invocations do add overhead. Moreover, our method is less effective on purely linear constraints, where concretization provides minimal benefit.
While \tool demonstrates notable solving time reductions on structured constraints, its robustness varies due to the introduction of \RL and \LLM. 
First, the effectiveness of \RL, typically \RL agents, heavily depends on the quality of the training set. 
Second, \tool under-perform on out-of-distribution constraints, which may cause \tool consistently select incorrect variables. 
Third, the performance of \LLM{} can vary. In this paper, we have selected the most effective \LLM by evaluating three popular \LLM{}s. However, it may introduce extensive effort to select the best one among all existing \LLM{}s. 
Finally, \tool is less effective on purely linear constraints, where concretization provides minimal benefit. This limitation stems from the fundamental nature of linear constraint solving and our method's design principles. Linear constraints, by definition, involve only first-degree polynomial relationships between variables, which modern SMT solvers like Z3 handle exceptionally well through specialized algorithms such as the simplex method and Fourier-Motzkin elimination~\cite{kroening2016decision}. These algorithms are highly optimized for linear arithmetic and can efficiently navigate the convex solution space without requiring the strategic variable concretization that \tool provides.

We observed this limitation empirically during our evaluation on the SMT-COMP \QFLIA{} benchmark, where \tool achieved a modest +4.6\% improvement compared to the more substantial +12.3\% improvement on \QFNIA{} (nonlinear) constraints. The underlying reason is that linear constraints lack the complex variable coupling and nonconvex search spaces that make strategic concretization beneficial. In linear systems, variables interact through simple additive relationships, and the constraint satisfaction problem reduces to finding points within well-defined polytopes. Concretizing variables in such systems often over-constrains the problem without providing the simplification benefits observed in nonlinear cases, where variable interactions involve products, powers, and modular arithmetic that create exponential search complexity.


\subsection{Threats to Validity}
Several potential threats to the validity of our findings should be considered when interpreting the experimental results.

\textbf{Internal Validity.} Our primary results are based on a specific configuration using LLaMA 3.1:70B as the LLM backend and Soft Actor-Critic (SAC) as the RL algorithm. While we conducted comparative evaluations with alternative LLM models (RQ2) and component ablations to validate design choices, the headline performance metrics reflect this primary configuration. The stochastic nature of both LLM generation and RL exploration introduces inherent variability in results. To mitigate these concerns, we employed several strategies: (i) consistent preprocessing pipelines across all experiments, (ii) solver-in-the-loop verification of every LLM-suggested assignment to ensure soundness, (iii) systematic ablation studies (RQ2) to isolate component contributions, and (iv) cross-solver validation (RQ1) to demonstrate architecture-agnostic benefits. All experiments were conducted under controlled computational environments with uniform timeout settings (1200 seconds, following SMT-COMP standards) to ensure reproducibility.

\textbf{External Validity.} Our evaluation encompasses two distinct benchmark categories with different theoretical foundations. The SMTimer dataset consists primarily of program-derived constraints involving bit-vector arithmetic from real-world software(Coreutils, BusyBox), while the SMT-COMP benchmarks provide \QFLIA{} and \QFNIA{} constraints focusing on linear and nonlinear integer arithmetic. This dual-dataset approach, while providing broader coverage than single-theory evaluations, still represents a limited subset of the full SMT theory spectrum. Key domains such as quantified formulas, arrays, strings, floating-point arithmetic, and uninterpreted functions remain largely unexplored beyond our specialized BVParti evaluation on bit-vector constraints.

Furthermore, our implementation relies on established open-source frameworks (Pearl for reinforcement learning and Ollama for LLM hosting), which introduces potential dependency-related validity concerns. While these frameworks provide robust, well-tested foundations that enhance reproducibility, they also constrain our experimental flexibility and may introduce framework-specific behaviors that could affect generalizability. The Pearl framework's SAC implementation and Ollama's model quantization strategies (Q4\_K\_M format) represent specific algorithmic choices that may not generalize to alternative RL algorithms or LLM deployment methods.

Additionally, our constraint filtering criteria (baseline solving time >300 seconds, >5 variables, non-\texttt{unsat} cases) may introduce selection bias toward specific problem characteristics, particularly favoring constraints with sufficient complexity to benefit from our approach while potentially excluding simpler cases where direct solving is more appropriate. The generalizability of our approach to other SMT theories, alternative constraint distributions, and different implementation frameworks requires further investigation.

\textbf{Construct Validity.} Our evaluation metrics—solving time and success rate—are inherently solver-centric and may not fully capture the complexity of constraint simplification effectiveness. These metrics are closely tied to the internal decision procedures and heuristics of specific solvers, potentially limiting the generalizability of performance comparisons. While we validated trends across multiple solver architectures (Z3, CVC5, MathSAT, BVParti) to address this concern, each solver's unique approach to handling concretized constraints may influence results differently. Furthermore, our hybrid reward system combines multiple signal sources (solver outcomes, predictive models, history-based penalties), and the relative weighting of these components may affect learning dynamics in ways not fully captured by end-to-end performance metrics.

\textbf{Conclusion Validity.} The choice of evaluation parameters can significantly influence result interpretation. Our adoption of the standard 1200-second timeout threshold, while consistent with SMT-COMP practices, may particularly impact constraints near this boundary where small performance shifts can change outcome classifications. Additionally, our selective simplification evaluation (RQ3) uses a specific routing threshold (time-bucket ≥4) that identifies approximately 0.8\% of constraints for \tool processing. This threshold was chosen for demonstration rather than optimization, and alternative thresholds may yield different cost-benefit trade-offs. The statistical significance of improvements, particularly for smaller datasets like BVParti (33 constraints), may be limited by sample size constraints.

These validity considerations underscore the importance of diverse experimental design, comprehensive cross-solver validation, and rigorous benchmarking when integrating machine learning approaches into symbolic reasoning systems. Future work should address these limitations through broader theory coverage, alternative evaluation metrics, and systematic sensitivity analyses across key experimental parameters.

% \subsection{Implications and Future Directions}
% Our results suggest a practical path toward solver-facing AI assistants. By reframing SMT simplification as \emph{policy learning}, \tool demonstrates that a machine-to-machine \emph{co-pilot}—a strategic \RL{} agent paired with a semantic \LLM{} and verified by the solver—can inject high-level guidance that brittle heuristics lack. Empirically, the gains generalize across heterogeneous solvers (RQ4), and a complexity-aware routing mechanism enables selective invocation under deployment budgets (RQ5), indicating that co-pilots are not only effective but also deployable.

% Building on this foundation, we outline several directions to strengthen the co-pilot both in capability and reliability:
% \begin{itemize}
%     \item \textbf{Safer, feedback-rich learning with experts in the loop:} Combine structure-aware signals (\eg, predicted complexity/impact, proof obligations) and text-to-reward guidance \cite{text2reward} with solver- and human-in-the-loop supervision; penalize unsound suggestions to stabilize training.
%     \item \textbf{Domain-specialized LLM and efficient deployment:} Fine-tune on SMT corpora and solver traces to improve semantic fidelity, then distill to lightweight policies and caching mechanisms to reduce latency and cost while preserving accuracy.
%     \item \textbf{Broader theories and transferable policies:} Extend to bit-vectors, arrays, uninterpreted functions, and quantified arithmetic (\eg, QF\_BV, AUFNIRA, UFNIA); design solver-invariant state/action abstractions that enable cross-solver transfer observed in RQ4.
%     \item \textbf{Budget- and risk-aware routing with trust guarantees:} Enhance routing with calibrated uncertainty and risk-sensitive objectives for dynamic, time-bucket-aware invocation (RQ5), and add post hoc verification and conformance checks to guard against distribution shift.
% \end{itemize}
% Ultimately, our work paves the way for a new class of hybrid reasoning systems where AI agents and symbolic engines collaborate to solve previously unsolvable problems, transforming the landscape of formal verification and program analysis.